\section{Architecture}

\subsection{Principes directeurs}

\begin{enumerate}
  \item \textbf{Logique métier pure dans \texttt{pomone-domain}} : zéro I/O, 100\,\% testable. Les types qui modélisent une plantation, un cycle, une récolte ne dépendent ni d'un SGBD ni de Slint.
  \item \textbf{Abstraction du SGBD via un trait \texttt{Repository}} dans \texttt{pomone-db} : SQLite et MariaDB partagent la même interface Rust. Les particularités SQL vivent dans les migrations et dans les implémentations concrètes, jamais dans la logique métier.
  \item \textbf{Pas de logique métier dans les triggers SQL} : tout calcul de dates, de rendements, de cycles est fait en Rust et testé. Les triggers Qrop ont été abandonnés volontairement pour éviter la duplication SQLite/MariaDB.
  \item \textbf{L'UI \texttt{pomone-ui} (Slint) ne contient pas de règles métier} : elle appelle uniquement des services exposés par \texttt{pomone-app}.
\end{enumerate}

\subsection{Crates}

\begin{description}
  \item[\texttt{pomone-domain}] Types métier purs (\texttt{Crop}, \texttt{Variety}, \texttt{Planting}, \texttt{Location}, \texttt{Strata}, \texttt{YearlyHarvest}\dots), validations, calculs de dates. Pas d'I/O, pas de dépendances I/O.
  \item[\texttt{pomone-db}] Trait \texttt{Repository} + implémentations concrètes (\texttt{sqlite/}, \texttt{mariadb/}), schémas et migrations (\texttt{migrations/sqlite/}, \texttt{migrations/mariadb/}), codec entre types domaine et types SQL.
  \item[\texttt{pomone-app}] Services applicatifs (\texttt{services.rs}), vues pour l'UI (\texttt{plantings\_view.rs}, \texttt{cultures\_view.rs}, \texttt{locations\_view.rs}), configuration, gestion d'état, i18n.
  \item[\texttt{pomone-ui}] Binaire desktop Slint (\texttt{pomone}). Écrans : Plantings, Cultures, Locations, Calendrier (en cours).
  \item[\texttt{pomone-cli}] Binaire admin / debug (\texttt{pomone-cli}). Utilitaires hors UI.
\end{description}

\subsection{Flux applicatif}

\begin{lstlisting}
   pomone-ui  (Slint)
       |
       v   appelle des services (use cases)
   pomone-app (services, etat applicatif, i18n)
       |
       v   Repository trait
   pomone-db  (sqlx — SQLite / MariaDB)
       |
       v   types
   pomone-domain (types purs, regles metier)
\end{lstlisting}

Le flux est strictement descendant : un écran Slint ne touche jamais à \texttt{sqlx} ni au schéma SQL ; un service \texttt{pomone-app} ne dépend que du trait \texttt{Repository} (pas de l'implémentation choisie au démarrage).
